\subsection{半单李代数的例子}

命题 8. 设 V 是有限维向量空间。则 \(\mathfrak{gl}(V)\) 是约化的；其中心是 V 的数乘变换所成的集合，其导出代数是 \(\mathfrak{sl}(V)\)，后者是半单的。

\(\mathfrak{gl}(\mathsf{V})\) 的恒等表示是单的，所以 \(\mathfrak{gl}(\mathsf{V})\) 是约化的，从而 \(\mathfrak{gl}(\mathsf{V})\) 是其中心 \(c\) 与其导出代数 \(\mathcal{D}(\mathfrak{gl}(\mathsf{V}))\) 的直和。中心 \(c\) 是数乘变换的集合（《代数》II 章§2，第 5 号，命题 5 的推论 1）。显然 \(\mathcal{D}(\mathfrak{gl}(\mathsf{V})) \subset \mathfrak{sl}(\mathsf{V})\)。由于 \(\mathfrak{sl}(\mathsf{V}) \cap c = \{0\}\)，\(\mathcal{D}(\mathfrak{gl}(\mathsf{V})) = \mathfrak{sl}(\mathsf{V})\)。因此 \(\mathfrak{sl}(\mathsf{V})\) 是半单的。

例。我们把 \(\mathfrak{sl}(\mathbf{K}^{2})\) 与二阶迹为零矩阵的李代数等同。记

\[
X = \left( \begin{array}{c c} 0 & 1 \\ 0 & 0 \end{array} \right) \qquad Y = \left( \begin{array}{c c} 0 & 0 \\ 1 & 0 \end{array} \right) \qquad H = \left( \begin{array}{c c} 1 & 0 \\ 0 & - 1 \end{array} \right).
\]

于是 \(X, Y, H\) 构成 \(\mathfrak{sl}(\mathbf{K}^2)\) 的一个基，并且

\[
[ H, X ] = 2 X \quad [ H, Y ] = - 2 Y \quad [ X, Y ] = H.
\]

由于维数为 1 或 2 的代数不是半单的（第 1 号，备注 1），\(\mathfrak{sl}(\mathbf{K}^2)\) 是单的。事实上，\(\mathfrak{sl}(\mathbf{V})\) 在 \(\dim V \geqslant 2\) 时是单的，我们稍后将看到这一点（另见习题 21 和 24）。

命题 9. 设 \(\mathrm{V}\) 为维数 \(n\) 的、定义在 \(\mathbf{K}\) 上的有限维向量空间，\(\beta\) 是 \(\mathrm{V}\) 上的非退化对称（分别为交错）双线性形式。令 \(\mathfrak{g}\) 为由 \(x \in \mathfrak{gl}(\mathrm{V})\) 组成的李代数，满足 \(\beta (xm, m') + \beta (m, xm') = 0\) 对所有 \(m, m'\)（属于 \(\mathrm{V}\)）成立。则 \(\mathfrak{g}\) 是约化的；\(\mathfrak{g}\) 是半单的，除非 \(\beta\) 对称且 \(n = 2\)。

对所有 \(u \in \mathfrak{gl}(\mathrm{V})\)，令 \(u^*\) 表示关于 \(\beta\) 的伴随；由《代数》IX 章§ 1，第 8 号的命题 7，有 \(\operatorname{Tr}(u) = \operatorname{Tr}(u^*)\)。条件

\[
\beta (u m, m ^ {\prime}) + \beta (m, u m ^ {\prime}) = 0
\]

对 V 中所有 \(m\)、\(m'\)，\(u + u^{*} = 0\) 成立。特别地，若 \(v \in \mathfrak{gl}(V)\)，则 \((v - v^{*})^{*} = v^{*} - v\)，因而 \(v - v^{*} \in \mathfrak{g}\)。现在令 \(u\) 为 \(\mathfrak{g}\) 中与 \(\mathfrak{g}\) 正交的元素，其中所用的双线性形式 \(\phi\) 与 \(\mathfrak{g}\) 的恒等表示相关。对所有 \(v \in \mathfrak{gl}(V)\)，\(\operatorname{Tr} u(v - v^{*}) = 0\) = 0，所以

\[
\operatorname{Tr} (u v) = \operatorname{Tr} (u v ^ {*}) = \operatorname{Tr} (u v ^ {*}) ^ {*} = \operatorname{Tr} (v u ^ {*}) = - \operatorname{Tr} (v u) = - \operatorname{Tr} (u v)
\]

从而 \(\operatorname{Tr}(u v) = 0\)。这说明 \(u = 0\)，所以 \(\phi\) 非退化。因此 \(\mathfrak{g}\) 是约化的（命题 5）。还需证明 \(\mathfrak{g}\) 的中心为零（除 \(\beta\) 对称且 \(n = 2\) 的情形外）。扩张基域后，可以假设 \(\mathbf{K}\) 是代数闭的。

\begin{enumerate}
\def\labelenumi{(\alph{enumi})}
\item
  当 \(\beta\) 对称时，它可以与 \(\mathbf{K}^n\) 上相对于典范基、矩阵为 \(I_{\mathrm{n}}\) 的双线性形式等同（《代数》IX 章§6，定理 1 的推论 1）。在这些条件下 \(\mathfrak{g}\) 可与斜对称矩阵的李代数等同（§3，第 4 号，例 1）。令 \(U = (u_{ij}) \in \mathfrak{g}\)；我们利用以下事实：\(U\) 与矩阵 \((v_{ij}) \in \mathfrak{g}\) 可交换，该矩阵除 \(v_{i_0j_0}\) 和 \(v_{j_0i_0}\) 外所有元素均为零（\(i_0 \neq j_0\)），且这两个元素分别等于 1 和 -1。我们得到 \(u_{i_0j} = u_{j_0j} = u_{ii_0} = u_{ij_0} = 0\)，当 \(i \neq i_0, j_0\) 且 \(j \neq i_0, j_0\) 时。若 \(n > 2\)，则对所有不同的指标 \(i_0\) 和 \(j\)，存在不同的指标 \(i\) 和 \(j_0\)，使得 \(i \neq i_0, j_0 \neq j, j_0 \neq i_0\)；所以 \(u_{i_0j} = 0\)。这证明 \(\mathfrak{g}\) 的中心元素为零。
\item
  当 \(\beta\) 交错且 \(n = 2m\) 时，\(\beta\) 可以与 \(\mathbf{K}^{2m}\) 上矩阵为 \(\left( \begin{array}{cc}0 & I_m \\ -I_m & 0 \end{array} \right)\) 的双线性形式等同（《代数》IX 章§5，定理 1 的推论 1）。在这些条件下 \(\mathfrak{g}\) 可与形如 \(U = \left( \begin{array}{cc}A & B\\ C & D \end{array} \right)\) 的矩阵李代数等同，其中 \(D = -^t A\)，\(B\) 和 \(C\) 是对称的（\(A, B, C, D\) 属于 \(\mathbf{M}_m(\mathbf{K})\)）（§3，第 4 号，例 1）。首先利用 \(U\) 与矩阵 \(\left( \begin{array}{cc}X & 0 \\ 0 & -^t X \end{array} \right)\) 可交换这一事实，其中 \(X \in \mathbf{M}_m(\mathbf{K})\)。于是 \(AX = X A\)、\(CX = -^t XC\)、\(XB = -B\)。\(^t X\)；由于这些等式必须对所有 \(X\) 成立，因而 \(A\) 是标量矩阵 \(\lambda I_m\)。现在利用 \(U\) 与矩阵 \(\left( \begin{array}{cc}0 & Y \\ 0 & 0 \end{array} \right)\) 可交换这一事实，其中
\end{enumerate}

\(Y\) 是 \(\mathbf{M}_m(\mathrm{K})\) 的对称矩阵。于是 \(\lambda Y = YC = CY = 0\)。这首先证明 \(\lambda = 0\)。此外，对所有 \(X \in \mathbf{M}_m(\mathrm{K})\)，\(X + {}^t X\) 是对称的，因而 \(XC = -{}^t XC\)。利用上面得到的等式 \(CX = -{}^t XC\)，可见 \(C\) 与 \(\mathbf{M}_m(\mathrm{K})\) 的每个元素可交换，所以 \(C\) 是标量矩阵；由于 \(YC = 0\)，它必为零。类似地可证明 \(B = 0\)。

当 \(\beta\) 对称且 \(n = 2\) 时，\(\mathfrak{g}\) 是一维的，因而是交换的。其他情形见习题 25 和 26。

